\section{The compiled polynomials with forty-one and forty-three coordinates}
\label{sec:p3-compiled-polynomials}

We apply the guarded finite-algebra theorem to the direct
Lucas-polynomial source of Section~\ref{sec:direct-integer-relation}.
All formal generators, monic relations, coefficient conversions,
denominator guards, and finite index ranges are supplied in the
corresponding polynomial specification.

\begin{theorem}\label{thm:p3-polynomial41}
Under the stated source arithmetic inputs and Sun's simultaneous
prime-pattern theorem, the supplied polynomial $P_{41}$ defines
$\mathbb Z$ in $\mathbb Q$ using forty-one total coordinates,
including the input. It has forty existential witnesses and satisfies
\[
        \deg P_{41}\le 2^{9566673},
        \qquad H(P_{41})\le 2^{\,2^{86098382}}.
\]
The semilocal-anchor alternative $P_{43}^{\mathrm{loc}}$ uses
forty-three total coordinates, forty-two existential witnesses, and
satisfies
\[
        \deg P_{43}^{\mathrm{loc}}\le 2^{236},
        \qquad H(P_{43}^{\mathrm{loc}})\le 2^{\,2^{240}}.
\]
Both are explicit indexed integer polynomials. These are upper bounds
for degree and height, not exact expansion statistics.
\end{theorem}
\begin{proof}
The source has forty retained rational coordinates and nineteen
formal algebra generators, nine of degree four and ten of degree
two. Its maximum finite rank is $2^{28}$. The semilocal alternative
adds two retained coordinates and one quadratic generator, giving
maximum rank $2^{29}$. Their rational input projections are
$\mathbb Z$ by Proposition~\ref{prop:direct-Lucas-preservation}.

The source-to-integer conversion is the same exact procedure proved
in Section~\ref{sec:compiled-shared-polynomials}. For every rational
square equation $q^2=N/D$, its denominator is guarded and the
substitution $Z=Dq$ gives $Z^2=ND$. Subsequent powers-of-two
clearings preserve the triangular monic presentation. Every
quadratic generator is split into its zero case and nonzero case,
with the appropriate closed equation or derivative guard. The nine
point quartics retain their derivative guards. Thus the required
geometric simplicity holds at every guarded point, including
exceptional parameter specializations.

Theorem~\ref{thm:explicit-guarded-finite-algebra} eliminates these
formal generators with one additional coordinate. It gives the
stated totals and preserves the rational projection of each source.
The emitted specifications have respectively 6018 and 6039 integer
coefficient-DAG nodes. The independent checker verifies the source
hash, all affine and rational-square changes, guard coverage, monicity,
triangularity, every finite index range, and the recursive degree and
coefficient majorants. The bounds are exactly those displayed above.
The fixed elliptic exclusion product is treated by the canonical
clearing and height estimate already proved for that product.
\end{proof}

The packages are
\nolinkurl{research/phase3/finite_fibre/polynomial-41/} and
\nolinkurl{research/phase3/finite_fibre/polynomial-43-semilocal/}.
Each contains its complete indexed specification, construction
report, and passing independent verification receipt. The large
multiplication matrices and outer products are not expanded.
The semilocal variant has a materialized coefficient DAG; the
forty-one-coordinate variant additionally contains the explicitly
indexed elliptic exclusion product.

The same producer also emits the two versions with input shift
$C=2402$ from Proposition~\ref{p3:integer-shift}. Their retained
counts, maximum ranks, and displayed degree and height bounds
are the same. They have 6019 and 6040 coefficient-DAG nodes,
respectively, and separate passing receipts. The additional
coefficient node is a fixed integer, not a quantified coordinate.
These variants remove the scalar-contact obstruction over fields
with a real embedding; they do not resolve the other arithmetic
obstructions over extensions.

\section{Quartics obtained from the direct Lucas test}
\label{sec:p3-quartics}

The direct test also improves the degree-four construction.
For this purpose we retain the scalar block with its preparation
coordinate, because its quadratic evaluation is less expensive
than the normalized scalar source used for finite-algebra
compression. We retain the proved comparison order $52$,
scalar order $32$, and fixed size constant
$\kappa=3^{2^{28}}$.

Replace the three proxy comparisons and the exact Lucas equation
by the one field-valued test
\[
 F(L_2,L_3,L_4),\qquad
 F(U,V,W)=V^2-U(U-2)VW+(U-2)^2W^2-1,
\]
of order $52$. Keep $L_1-\beta$ of order $16$ and all four
two-anchor cross comparisons. The direct-polynomial congruence
lemma applies with $k=52$ just as with $k=64$.
The precision-$52$ height proof uses the same integer residual
$\Delta$ and its bound $|\Delta|<5B^4$, so its threshold and
the chosen anchor are unchanged. The scalar-order-$32$ proof
is unaffected. Completeness uses the same fixed-point
neighborhoods: a polynomial in the original logarithm ratios
remains bounded there, introduces no new denominator, and
has the required prime-power precision.

The quadratic compiler evaluates the new field test by computing
\[
 A=U(U-2),\qquad B=VW,\qquad C=(U-2)W,
 \qquad F=V^2-AB+C^2-1.
\]
Each field multiplication uses its two rational component equations
as before. The coordinate detectors now have ten tests of order
one, two of order sixteen, and four of order fifty-two.
Their symmetric-product coefficient lists are shared across all
nine coordinate detectors. This removes more arithmetic
coordinates than merely deleting the three proxies would do.
All reciprocal-domain lemmas, Boolean scalar choices, and
quadratic-preserving substitutions from the preceding quartic
construction are retained.

\begin{theorem}\label{thm:p3-quartic1315}
Under the arithmetic inputs stated for the shared construction,
the supplied polynomial $P_{1315}$ defines $\mathbb Z$ in
$\mathbb Q$, has 1315 total variables and degree exactly four.
It is the sum of squares of 1272 primitive quadratic residuals
and has 28,991 nonzero encoded entries in its complete table over
$\mathbb Z[\kappa]$, with the specified fixed integer
$\kappa=3^{2^{28}}$. After collecting powers of this constant,
there are 28,987 distinct monomials in the quantified coordinates.

The alternative $P_{1339}$ has 1339 total variables, degree
exactly four, 1296 quadratic residuals, and 29,061 expanded
ordinary integer terms. Its primitive coefficient height is
less than $2^{1020}$.
\end{theorem}
\begin{proof}
The preceding argument proves the arithmetic source replacement
at precision $52$. The exact quadratic compilation and the
domain equivalences were proved in Section~\ref{sec:new-quartic}.
For the present source the independent checker verifies 1285
arithmetic gates, eleven row additions, and sixty coordinate
eliminations. It then independently squares and convolves every
residual and compares the complete resulting coefficient table.
The active coordinate count is 1315, the number of primitive
residuals is 1272, and the number of nonzero table entries is
28,991. The encoding separates different powers of $\kappa$
in a coefficient. Collecting these powers merges four pairs of
entries: the constant coefficient is positive, of the form
$c_0+1024\kappa^2$, and three further coefficients are
$1-192\kappa\ne0$. Thus there are 28,987 monomials in the
actual coordinates, with no cancellation after substituting the
specified integer. A nonzero degree-four coefficient occurs, so the degree
is exactly four. Over $\mathbb Q$ a sum of squares vanishes
exactly when all its residuals vanish, giving the asserted
input projection.

For the ordinary integer alternative, retain the same twenty-four
successive squaring coordinates computing $\kappa$ from $3^{16}$.
The count becomes 1339, and the checker verifies 1309 arithmetic
gates with the same eleven row additions and sixty eliminations.
Its 1296 primitive residuals give exactly 29,061 nonzero integer
terms. Their largest coefficient has bit length at most 1020,
and their content is one. This proves the degree and height
claims without leaving any coefficient expansion implicit.
\end{proof}

The two tables are in
\nolinkurl{research/phase3/quartic/quartic-direct/} and
\nolinkurl{research/phase3/quartic/quartic-direct-small-coefficients/}.
Their producer is \nolinkurl{build_direct_quartic.py}.
The independently written checker
\nolinkurl{research/phase2/quartic/verify_shared_quartic.py}
is reused without modification and does not import this producer.
These algebraic checks do not certify the external arithmetic
lemmas. The number-field section separately examines both
literal quartics and faithful combinations of their residual systems.
